% !TeX root = ../Thesis.tex
\chapter{苏格拉底式对话教育智能体的设计框架}
\label{Detailed_Design_Framework}

面向初中物理迷思概念干预的苏格拉底式对话教育智能体，核心并非单纯调用大语言模型生成自然语言回复，而是将教师在概念干预过程中的启发式提问、逻辑追问与支架支持，转化为可执行、可控制、可评估的系统机制。据此，本研究围绕“教学策略算法化建模—对话控制机制设计—领域知识支撑”三个层面，构建苏格拉底式对话教育智能体的整体设计框架。该框架可保障智能体在多轮对话中维持教学逻辑的一致性，同时针对学生的典型迷思概念开展精准启发式引导，最终实现“引导理解”而非“直接代答”的教育目标。


\section{教学策略的算法化建模}

苏格拉底教学法以追问、反思、澄清、论证为核心手段，引导学习者自主发现认知偏差与逻辑矛盾，进而完成科学概念的重建。对教育智能体而言，这一教学过程无法仅通过单一提示词或开放式语言生成实现，需将依赖教师教学经验的主观教学行为转化为结构化、可计算的对话策略规则，把隐性的启发式教学策略编码为计算机可识别、可执行、可管控的标准化算法逻辑，以解决通用Prompt驱动下智能体教学行为不可控、教学逻辑断裂的问题。

据此，本研究并非将苏格拉底式教学简单理解为“连续发问”，而是将其视为一个包含“反讽—助产—归纳—定义”的递进式教学过程：先通过反诘暴露学生原有认知中的内在矛盾，再通过助产式引导支持学生生成新的解释，最终借助归纳与定义推动其形成较为稳定的科学概念表征。在此基础上，本研究进一步引入Paul-Elder批判性思维框架，将学生回答中涉及的概念表述、隐含假设、证据基础、推理后果等维度纳入提问分析范围，并结合波斯纳概念转变理论中“引发不满—新概念可理解—新概念具合理性—新概念富有成效”的四个条件，对初中物理迷思概念干预中的提问行为开展算法化建模，构建面向迷思概念干预的分层提问策略体系。

\subsection{适配迷思概念转变的分层提问策略库}

本研究融合苏格拉底“反讽—助产—归纳—定义”的核心教学流程、Paul-Elder批判性思维框架与波斯纳概念转变理论，针对初中物理迷思概念的认知特征，构建包含澄清提问、挑战假设、证据追问、后果探索、类比支架的五类核心提问策略体系。需要说明的是，前四类策略主要对应Paul-Elder框架中的澄清、假设、证据与后果维度，第五类“类比支架”则是在此基础上，结合科学教育中类比支架与支架式教学研究，为适配抽象物理概念理解需求所作的教育化扩展。各类策略对应明确的教学目标与适用场景，功能上互不割裂，可在对话过程中根据学生的应答状态动态组合、逐层递进，具体内容见表~\ref{tab:socratic_strategies}。

上述策略共同构成教育智能体的核心提问策略体系。其本质并非固定话术模板的简单堆叠，而是适配不同迷思概念场景的教学行为规则集合。系统运行过程中，将根据学生的实时应答内容、迷思类型与认知状态，动态匹配对应策略，保障苏格拉底式对话的递进性与针对性。需要进一步说明的是，Paul-Elder框架中的“观点与视角”维度在本研究中未单列为独立策略类型，而是吸纳于挑战假设、后果探索与类比支架的使用过程中，用于引导学生从不同情境、不同模型或不同解释框架审视原有判断，以保持策略体系的简洁性与教学适配性。

\begin{table}[htbp]
\centering
\caption{苏格拉底式对话核心提问策略库}
\label{tab:socratic_strategies}
\resizebox{\textwidth}{!}{%
\begin{tabular}{|p{2.8cm}|p{3.6cm}|p{2.5cm}|p{2.8cm}|p{3.2cm}|p{3.8cm}|p{6.2cm}|}
\hline
\textbf{策略类型} & \textbf{核心教学目标} & \textbf{主要理论依据} & \textbf{对应Paul-Elder维度} & \textbf{对应概念转变条件/环节} & \textbf{触发条件} & \textbf{算法动作逻辑与物理场景示例} \\
\hline

澄清提问（Clarification） 
& 消除语义歧义，引导学生完成从日常模糊表达到精准科学表达的转化，暴露隐含的认知偏差 
& 苏格拉底式澄清追问；Paul-Elder提问分类中的“澄清问题” 
& 清晰性（Clarity）、精确性（Precision）、概念（Concepts） 
& 新概念可理解的前置条件；迷思概念识别阶段 
& 1. 学生输入包含模糊性、非科学的口语化表述（如“电没了”“东西变轻了”）；2. 学生对物理量的表述混淆（如电压/电流不分、质量/重力混淆） 
& 算法动作：提取输入中的模糊词汇，结合当前对话的知识点场景，生成封闭式精准追问，引导学生明确表述的科学内涵。示例：学生提到“灯泡让电变少了”，生成追问：“你提到的‘电变少了’，具体是指电流减小、电荷减少，还是电能减少了？” \\
\hline

挑战假设（Assumption Probing） 
& 揭示学生迷思概念背后的隐含前提与直觉错误，制造认知冲突，打破学生对原有错误概念的认知自洽 
& 苏格拉底反诘法；Paul-Elder提问分类中的“探查假设” 
& 假设（Assumptions）、逻辑性（Logic） 
& 引发不满；认知冲突制造阶段（反讽环节） 
& 1. 学生基于日常直觉做出错误判断，明确表达了迷思概念对应的核心假设（如“电池没电是因为电荷被灯泡用光了”）；2. 学生的结论与基本物理定律相悖 
& 算法动作：锁定学生表述中的核心错误假设，调用知识库中的物理定律、常识反例，生成指向性反问，引导学生自主发现原有认知的内在矛盾。示例：学生提出“电荷被灯泡消耗了”，生成反问：“如果电荷在流经灯泡时被用光了，那电荷这种物质就凭空消失了，这是否违背了我们学过的电荷守恒定律？” \\
\hline

证据追问（Evidence Seeking） 
& 培养学生的实证意识与科学推理能力，引导学生完成从“主观猜想”到“基于证据的科学论证”的转变，强化科学探究思维 
& 苏格拉底论证追问；Paul-Elder提问分类中的“探查理由与证据” 
& 信息/证据（Information/Evidence）、准确性（Accuracy）、相关性（Relevance） 
& 新概念具合理性；概念反思阶段 
& 1. 学生提出主观猜想、判断，但未给出对应的实验现象或推理依据；2. 学生仅凭直觉下结论，未关联物理原理与证据 
& 算法动作：围绕学生的猜想，生成指向证据来源、推理逻辑的追问，引导学生回溯结论的推导过程。示例：学生说“重的物体在水里沉得更快”，生成追问：“你是通过什么实验现象得出这个结论的？能不能用我们学过的浮力公式来解释这个判断？” \\
\hline

后果探索（Consequence Exploration） 
& 通过归谬法推演错误逻辑的衍生结论，进一步强化认知冲突，让学生自主发现迷思概念的逻辑漏洞，完成概念解构 
& 苏格拉底归谬式反问；Paul-Elder提问分类中的“探查后果与影响” 
& 后果与影响（Implications and Consequences）、逻辑性（Logic）、广度（Breadth） 
& 引发不满的深化；认知冲突深化阶段（反讽环节） 
& 1. 学生坚持错误的物理模型，对单次挑战无明显认知动摇；2. 学生的错误逻辑会衍生出与生活常识、物理定律相悖的推论 
& 算法动作：基于学生的错误模型进行逻辑推演，生成极端化、生活化的归谬反问，引导学生直观感知原有逻辑的谬误。示例：学生坚持“串联电路里电流被灯泡逐个消耗”，生成反问：“按照这个逻辑，家里串联的灯泡越多，后面的灯泡会越暗，那你家开的灯越多，邻居家的灯也会变暗，这符合你的生活经验吗？” \\
\hline

类比支架（Analogical Scaffolding） 
& 借助学生熟悉的生活化场景搭建认知脚手架，降低抽象物理概念的理解门槛，引导学生完成从错误模型到科学模型的自主建构 
& 支架式教学；类比教学；科学教育中的类比支架研究 
& 概念（Concepts）、视角/观点（Point of View）、相关性（Relevance） 
& 新概念可理解；新概念建构阶段（助产环节） 
& 1. 学生经历认知冲突后陷入认知僵局，无法自主构建正确的科学模型；2. 学生对抽象的物理概念存在理解障碍；3. 检测到学生出现挫败、畏难的情绪状态 
& 算法动作：检索知识库中对应知识点的标准化教学类比案例，生成贴合初中生认知的类比引导，不直接给出概念定义，通过类比引导学生自主推导。示例：学生无法理解电压、电流的关系，生成引导：“我们可以把电路想象成封闭的水管回路，电压好比水管两端的水压，电流好比水管里的水流，你可以试着用这个类比，说说电阻会对应水管里的什么？” \\
\hline

\end{tabular}%
}
\end{table}

\subsection{教育场景中的禁泄露约束机制}

教育场景中，生成式人工智能应用的核心风险之一是“过度帮助”，即在学生未完成必要的思考过程前，直接输出答案、公式或结论，削弱学生自主建构科学概念的机会。据此，本研究将“禁泄露”机制定位为苏格拉底式对话教育智能体的核心约束条件，而非附加功能模块。

本研究从两个维度构建禁泄露机制：一是系统级教学约束，即在智能体的基础角色设定中，明确其教育身份与行为边界——系统并非答题工具，而是引导学生自主思考的教学支持者，核心任务并非直接输出结论，而是通过追问、支架搭建与过程验证，推动学生自主形成对概念的科学理解。二是输出审核约束，即在系统生成候选回复后，开展二次合规性审查，识别内容中是否包含直接结论、完整标准答案、一步到位的解题路径，或其他可替代学生思考过程的内容；一旦检测到高风险回复，将触发重生成指令，要求系统输出更具启发性的引导内容。

需要明确的是，本研究设计的禁泄露机制并非简单的“拒绝回答”，而是在拒绝直接代答的同时，为学生提供适配的思维引导支持。其设计目标并非削弱系统的教学支持能力，而是在教育有效性与教学安全性之间实现平衡，保障智能体始终恪守“引导而非代答”的核心教学原则。

\section{Agent核心架构设计}

为保障苏格拉底式教学策略的稳定落地，本研究采用“感知—决策—执行”三层架构，对智能体的认知与交互过程进行解耦设计。该架构可实现学生输入分析、教学状态控制与自然语言生成的分模块处理，提升系统的可解释性与可维护性，同时可基于LangGraph框架便捷实现对话状态流转、循环控制与上下文管理。

\subsection{感知层：学生输入分析与认知状态识别}

感知层的核心任务，是对学生输入开展教学导向的语义解析，而非通用的自然语言理解。即系统需识别的不仅是“学生说了什么”，更核心的是学生当前的认知状态、是否暴露特定迷思概念、是否存在索要答案的意图或负面学习情绪等与教学决策直接相关的关键信息。

基于上述目标，本研究将感知层的分析输出构建为四个维度的标签体系，具体见表~\ref{tab:perception_tags}。基于感知层输出的多维标签，系统可完成对学生当前认知状态的初步表征，为后续的教学状态判定与策略调用提供核心依据。

\begin{table}[htbp]
\centering
\caption{感知层多维度标签体系}
\label{tab:perception_tags}
\resizebox{\textwidth}{!}{%
\begin{tabular}{|p{3cm}|p{10cm}|p{4cm}|}
\hline
\textbf{标签类型} & \textbf{标签定义与分类} & \textbf{功能说明} \\
\hline
对话意图标签（Intent） & 基于预设分类体系，识别学生输入的核心意图，共分为6类：Direct\_Answer\_Seek（直接索要答案）、Misconception\_Expression（表达迷思概念）、Hypothesis\_Put\_Forward（提出猜想/假设）、Cognitive\_Stuck（表达认知困惑/陷入僵局）、Knowledge\_Inquiry（概念询问）、Off\_Topic（偏离教学主题的闲聊） & 为后续状态跳转提供核心依据 \\
\hline
迷思概念标签（Misconception\_Tag） & 基于本研究构建的初中物理迷思概念库，通过语义匹配识别学生输入中隐含的迷思概念，输出对应迷思ID、所属知识点与干预要点，例如识别出学生持有“电流消耗模型”“浮力与深度正相关”等典型迷思 & 直接对接知识库中的对应干预策略 \\
\hline
认知状态标签（Cognitive\_State） & 结合对话历史与当前输入，判定学生所处的概念转变阶段，共分为5个阶段：原有概念固守、认知冲突触发、认知僵局、新概念探索、概念掌握验证 & 为教学策略的选择提供依据 \\
\hline
情感状态标签（Sentiment） & 识别学生的学习情绪状态，共分为4类：Frustrated（挫败/畏难）、Curious（好奇/积极）、Confused（困惑/迷茫）、Impatient（急躁/不耐烦） & 检测到负面情绪时，决策层将调整策略的激进程度，优先增加支架支持，避免加剧学生的学习挫败感 \\
\hline
\end{tabular}%
}
\end{table}

\subsection{决策层：基于教学状态机的对话规划}

决策层是智能体的核心控制单元，核心任务是基于感知层的分析结果、对话上下文与预设教学目标，确定系统下一步的教学行动。为避免对话过程退化为无状态、线性的问答生成，本研究基于LangGraph设计教学状态机，对苏格拉底式引导的全流程进行显式建模，具体内容见表~\ref{tab:teaching_state_machine}。

\begin{table}[htbp]
\centering
\caption{教学状态机核心状态与流转规则}
\label{tab:teaching_state_machine}
\resizebox{\textwidth}{!}{%
\begin{tabular}{|p{2cm}|p{6cm}|p{8cm}|}
\hline
\textbf{状态ID} & \textbf{状态名称与定位} & \textbf{处理逻辑与跳转规则} \\
\hline
S0 & 监听与分析状态（Listen \& Analyze）\newline【对话入口默认状态】 & 接收学生输入文本，调用感知层完成多维标签的识别与输出，同步更新全局对话状态，处理完成后根据标签结果触发状态跳转。 \\
\hline
S1 & 护栏检查与合规处理状态（Guardrails Check）\newline【前置安全校验状态】 & 接收S0的输出结果，核心校验两类风险：一是学生输入意图为“直接索要答案”，二是输入内容偏离教学主题或不符合教育规范。若触发风险，直接跳转至S2拒绝与引导状态；若合规无风险，跳转至S3迷思诊断状态。 \\
\hline
S2 & 拒绝与引导状态（Refusal \& Guidance） & 针对学生直接索要答案、代写作业等请求，执行标准化的拒绝与引导流程：既明确拒绝直接代答的要求，又通过引导性提问将学生拉回自主思考的路径，避免生硬拒绝降低学生的学习体验。处理完成后，跳转回S0等待下一轮学生输入。 \\
\hline
S3 & 迷思概念诊断状态（Misconception Diagnosis） & 基于感知层输出的迷思标签、认知状态标签，判定学生是否持有明确的迷思概念，以及其所处的概念转变阶段。若检测到明确的迷思概念，跳转至S4反讽与认知冲突状态；若未检测到明确迷思，学生处于概念探索/困惑状态，跳转至S5助产与支架引导状态；若学生已完成概念修正、给出正确表述，跳转至S6验证与深化状态。 \\
\hline
S4 & 反讽与认知冲突状态（Elenchus \& Cognitive Conflict）\newline【对应苏格拉底“反讽”环节】 & 核心目标是打破学生对原有迷思概念的认知自洽。该状态下，决策层将匹配对应迷思概念的干预要点，选择“挑战假设”“后果探索”策略，向执行层下发策略生成指令，引导学生自主发现原有认知的内在矛盾。处理完成后，跳转回S0，基于学生的下一轮输入重新判断状态。 \\
\hline
S5 & 助产与支架引导状态（Maieutics \& Scaffolding）\newline【对应苏格拉底“助产术”环节】 & 核心目标是为陷入认知僵局的学生搭建认知脚手架，引导学生自主建构正确的科学模型。该状态下，决策层将根据学生的认知状态、情绪状态，选择“澄清提问”“证据追问”“类比支架”策略，向执行层下发生成指令，逐步引导学生完成新概念的建构。处理完成后，跳转回S0，基于学生反馈更新状态。 \\
\hline
S6 & 验证与深化状态（Verification \& Deepening）\newline【对应苏格拉底“归纳-定义”环节】 & 核心目标是验证学生是否真正掌握科学概念，避免出现“蒙对答案但未理解概念本质”的问题，同时实现知识的迁移与深化。该状态下，决策层将生成追问，要求学生解释推理过程、说明原理依据，或提出变式问题检验概念的迁移能力。完成验证后，跳转回S0，等待新一轮对话输入。 \\
\hline
\end{tabular}%
}
\end{table}

教学状态机具有以下核心特性：

一是全局状态持久化。基于LangGraph的Checkpointer机制，实现对话全流程的状态持久化存储，完整记录学生的迷思概念表达、认知状态变化与历史对话逻辑，保障智能体在多轮对话中可稳定追踪学生的核心认知偏差，避免出现上下文遗忘、逻辑前后矛盾的问题。

二是自适应状态跳转与回溯。状态机并非线性推进，而是根据学生的实时反馈动态调整。例如，学生在完成新概念建构后，若再次出现原有迷思概念，状态机将自动回溯至S3诊断状态与S4反讽状态，重新触发认知冲突干预，而非继续推进既定流程，可充分适配学生概念转变过程的反复性特征。

三是教学节奏动态管控。基于学生的情感状态与认知水平，动态调整教学策略的激进程度。例如，检测到学生出现挫败情绪时，可从S4反讽状态回退至S5支架引导状态，降低认知负荷，保障对话始终处于学生的最近发展区内。

该状态机设计使教学智能体的对话行为不再完全依赖大语言模型的即时生成倾向，而是受到显式教学逻辑的严格约束。系统每一次回复的生成，均需经过“当前状态—触发条件—策略选择—下一状态”的完整规划流程，可显著提升多轮对话的逻辑一致性、行为可控性与教学针对性。

\subsection{执行层：面向教学任务的受控语言生成}

执行层的核心功能，是将决策层输出的教学策略指令，转化为符合初中生认知水平的自然语言回复。与通用大语言模型的开放式生成不同，本研究的执行层采用“多约束动态Prompt组装”机制，保障输出内容严格服务于预设教学目标，避免出现偏离教学逻辑或过度代答的问题。

\[
\begin{split}
\text{最终Prompt} =& \underbrace{\text{核心教学身份指令}}_{\text{固定角色设定}} + \underbrace{\text{当前状态与策略执行要求}}_{\text{决策层输出}} \\
&+ \underbrace{\text{迷思概念匹配知识片段}}_{\text{知识库检索结果}} + \underbrace{\text{禁泄露约束规则}}_{\text{系统级安全要求}} \\
&+ \underbrace{\text{近三轮对话上下文摘要}}_{\text{语境连贯性保障}}
\end{split}
\]

该机制通过四类核心信息的整合实现生成过程的可控性：
\begin{enumerate}
    \item 教学目标锚定信息：包含当前对话所处的教学状态、选定的提问策略类型与具体执行要求，保障回复严格匹配苏格拉底式教学的阶段性目标，避免生成脱离教学逻辑的内容。
    \item 认知状态适配信息：包含学生当前的迷思概念标签、认知阶段与情感状态识别结果，保障回复的难度与引导方式适配学生的最近发展区。
    \item 学科知识支撑信息：包含从领域知识库中检索到的对应迷思概念的反例、类比素材与干预要点，从源头降低生成内容出现学科性错误的风险。
    \item 输出边界约束信息：包含禁泄露规则与教学行为边界限定，从生成层面直接阻断直接输出答案、完整解题步骤等可替代学生思考过程的内容。
\end{enumerate}

例如，当决策层判定需采用“后果探索”策略针对“电流消耗模型”开展干预时，执行层将基于上述框架自动组装Prompt，生成指向错误逻辑内在矛盾的归谬式提问，而非直接陈述“串联电路中电流处处相等”的科学结论，既保障语言的自然流畅，又严格恪守“引导而非代答”的教学原则。

执行层的设计本质，是将大语言模型的生成能力限定在教学策略预设的边界内，实现“教学意图-语言输出”的精准映射，在保障回复流畅性的同时，使其完全服务于迷思概念转变的核心教学目标。

\section{初中物理领域知识库构建}

本模块旨在解决大模型生成幻觉、引导内容不符合课标要求、迷思概念干预精准度不足等问题，为智能体的对话引导提供精准、权威、贴合本土教学实际的学科知识支撑，保障所有引导内容严格契合人教版初中物理教学要求与《义务教育物理课程标准（2022年版）》。

\subsection{知识库的内容体系构建}
本知识库严格围绕人教版初中物理八年级核心教学模块，聚焦本研究选定的电学（电流与电路、欧姆定律）、力学（浮力）两大迷思概念高发模块，构建结构化知识体系。知识库所有内容均来自权威数据源，包括《义务教育物理课程标准（2022年版）》、人教版初中物理教材、国内初中物理迷思概念领域的核心研究文献、一线骨干教师的优质教学案例。知识库分为三大核心子库。

核心概念知识子库梳理电学、浮力模块课标要求的核心知识点，明确每个知识点的科学定义、核心物理规律、公式内涵与教学重难点，为智能体的引导内容提供权威学科依据，避免出现知识错误与生成幻觉。

典型迷思概念子库是知识库的核心模块。针对电学、浮力模块，系统梳理经学术研究与教学实践验证的典型迷思概念，为每个迷思概念构建结构化信息卡片，内容涵盖迷思ID、所属知识点、迷思核心内涵、学生形成该迷思的认知根源、对应的苏格拉底干预策略、教学反例与类比案例，实现“迷思识别-干预策略的精准映射。

教学策略与类比案例子库收集整理初中物理一线教学中成熟的、符合初中生认知特点的生活化类比案例、探究式提问案例、概念转变引导脚本，为执行层的回复生成提供标准化教学参考，提升引导内容的教学适配性。

\subsection{核心模块的迷思概念结构化编码}
本研究对电学、浮力两大模块的典型迷思概念开展标准化编码，实现与感知层、决策层的精准联动，具体内容见表~\ref{tab:misconception_coding}。
\begin{table}[htbp]
\centering
\caption{初中物理核心模块迷思概念结构化编码表}
\label{tab:misconception_coding}
\resizebox{\textwidth}{!}{%
\begin{tabular}{|p{2cm}|p{6cm}|p{2cm}|p{7cm}|}
\hline
\textbf{所属模块} & \textbf{覆盖核心概念} & \textbf{迷思编码} & \textbf{迷思概念内涵} \\
\hline
\multirow{4}{=}{电学模块（电流与电路、欧姆定律）} & \multirow{4}{=}{电流、电压、电阻、电荷守恒定律、串联与并联电路的规律} & M-ELE-001 & 电流消耗模型：认为电流在流经用电器时会被消耗，串联电路中电流从电源正极流出后会逐渐减小 \\
\cline{3-4}
& & M-ELE-002 & 单极模型：认为电流只需要从电池正极流出就能点亮灯泡，不需要形成闭合回路 \\
\cline{3-4}
& & M-ELE-003 & 电压/电流混淆：将电压与电流的概念等同，认为“有电压就一定有电流”“灯泡不亮是因为没有电压” \\
\cline{3-4}
& & M-ELE-004 & 顺序推理谬误：认为开关必须接在灯泡“前面”（靠近电源正极一侧）才能控制灯泡，接在“后面”就无法控制 \\
\hline
\multirow{4}{=}{浮力模块} & \multirow{4}{=}{阿基米德原理、浮力的产生原因、物体的浮沉条件、液体密度与排开液体体积对浮力的影响} & M-BUO-001 & 重物下沉论：认为物体的浮沉完全由重量决定，重的物体一定会下沉，轻的物体一定会上浮，忽略物体密度与排开液体体积的影响 \\
\cline{3-4}
& & M-BUO-002 & 深度浮力正相关谬误：认为物体在水中的深度越深，受到的浮力越大，混淆了浮力与液体压强的概念 \\
\cline{3-4}
& & M-BUO-003 & 漂浮浮力不等论：认为漂浮在水面的不同物体，受到的浮力大小不同，忽略了漂浮时浮力等于物体重力的核心规律 \\
\cline{3-4}
& & M-BUO-004 & 孔洞下沉谬误：认为有孔洞、漏水的物体一定会下沉，忽略了物体整体密度与液体密度的关系 \\
\hline
\end{tabular}%
}
\end{table}

\subsection{基于RAG的知识检索与调用机制}
本研究采用ChromaDB向量数据库实现知识库的存储与检索，并将其嵌入智能体对话流程，具体调用逻辑如下：
\begin{enumerate}
\item 知识切片与向量化：将知识库中的结构化内容按照知识点和迷思概念进行细粒度切片，通过开源嵌入模型完成向量化处理后存储至向量数据库，同时保留结构化标签，以支持精准的条件检索。
\item 多触发点检索机制：在对话流程中设置三类检索触发节点，以提高知识供给的针对性。其一，当感知层识别出对应的迷思概念标签时，系统自动检索该迷思对应的干预要点、反例与类比案例，为决策层的策略选择提供依据；其二，当状态机跳转至相应教学状态时，系统检索对应知识点的课标要求与教学重难点，以保证引导内容符合教学规范；其三，执行层在组装Prompt时自动调用检索结果，确保生成内容与目标知识点保持一致，减少大模型幻觉。
\item 检索结果优化：采用“关键词精准匹配+语义相似度检索”的双路检索策略，并设置检索结果的置信度阈值，以保证召回内容的相关性与准确性，减少无关信息对模型生成的干扰。
\end{enumerate}
